\begin{textsample}{POS Dim 2 – 2011\_c6 – Score -2.00 – 2011\_c6/enade\_2011\_art\_vis\_14aea27cb0.txt}  \label{ex:f2_pos_410}
CAPÍTULO 3 ANÁLISE TÉCNICA DA PROVA Este capítulo tem por objetivo apresentar o desempenho dos estudantes concluintes de Licenciatura em Artes Visuais no ENADE/2011.
Para isso, foram calculadas as estatísticas básicas da prova em seu todo, bem como as estatísticas dos componentes relacionadas à Formação Geral, ao de Conhecimento Específico da Área e das questões discursivas isoladamente.
Nas tabelas, são apresentados o tamanho da população inscrita e de presentes, e as seguintes estatísticas das notas8 : média do desempenho na prova, erro padrão da média, desvio padrão, nota mínima, mediana e nota máxima.
As estatísticas apresentadas neste capítulo contemplam o total de estudantes concluintes da área de Licenciatura em Artes Visuais em 2011 do Brasil e, separadamente, por Grande Região.
Foram calculadas tendo -se em vista as seguintes agregações : ( a ) as Grandes Regiões e o país como um todo ; ( b ) a Categoria Administrativa ; e ( c ) a Organização Acadêmica.
Em relação aos gráficos de distribuição de notas, o intervalo considerado foi de 10 unidades, aberto à esquerda e fechado à direita, com exceção do primeiro intervalo, [ 0 ; 10 ], fechado em ambos os extremos.
Para os gráficos de distribuição das notas das questões discursivas, foram consideradas mais duas categorias : questão em branco e nota zero. 3.1 ESTATÍSTICAS BÁSICAS DA PROVA 3.1.1 Estatísticas Básicas Gerais A Tabela 3.1 apresenta as estatísticas básicas da prova por grande Região.
A população total de inscritos foi de 3.712.
Destes, 2.916 estiveram presentes, sendo 21,4 \% o índice de não comparecimento.
A Região de maior abstenção foi a Norte ( 41,1 \% ) e a de menor abstenção foi a Sul ( 9,8 \% ).
A média das notas da prova como um todo ( nas seções seguintes serão analisados os componentes de Formação Geral e de Conhecimento Específico ) foi 51,1, sendo que os alunos da região Sudeste obtiveram a média mais baixa ( 48,7 ) e os da região Nordeste obtiveram a média mais alta ( 57,8 ).
As demais médias foram : 54,5 na região Norte, 52,3 na região Sul e 53,6 na região Centro-Oeste.
O desvio padrão para o Brasil como um todo foi 15,7, sendo o maior desvio padrão encontrado na região Sudeste ( 17,3 ).
Nas demais regiões o desvio padrão de manteve entre 13,1 na região Centro-Oeste e 13,9 na região Sul, indicando haver sempre uma menor dispersão das notas do que na região Sudeste. 8 Essas estatísticas e outras estão definidas no Capítulo 1.
O comportamento das notas dos estudantes de todo o Brasil pode ser observado no Gráfico 3.1, que apresenta um histograma com a distribuição das mesmas.
Essa é uma distribuição unimodal com moda no intervalo ( 50 ; 60 ].
O coeficiente de assimetria da distribuição das notas é negativo ( 0,70 ), indicando um maior espalhamento das notas abaixo do intervalo modal e maior concentração acima dele.
As distribuições por Grande Região também apresentam assimetria negativa, maior em valor absoluto nas regiões Norte e Nordeste, 1,17 e 0,84, respectivamente. 15,7, sendo o maior desvio padrão encontrado na região Sudeste ( 17,3 ).
Nas demais regiões o desvio padrão de manteve entre 13,1 na região Centro-Oeste e 13,9 na região Sul, indicando haver sempre uma menor dispersão das notas do que na região Sudeste.
A região que obteve a maior nota máxima foi a Centro-Oeste ( 86,7 ), ao passo que a região que atingiu a menor nota máxima foi a Norte ( 81,7 ).
A mediana do Brasil como um todo foi 53,0, sendo a maior mediana obtida na região Nordeste ( 59,3 ) e a menor obtida na Sudeste ( 51,4 ).
A nota mínima só foi diferente de zero ( 21,1 ) na região Centro-Oeste.
O comportamento das notas dos estudantes de todo o Brasil pode ser observado no Gráfico 3.1, que apresenta um histograma com a distribuição das mesmas.
Essa é uma distribuição unimodal com moda no intervalo ( 50 ; 60 ].
O coeficiente de assimetria da distribuição das notas é negativo ( 0,70 ), indicando um maior espalhamento das notas abaixo do intervalo modal e maior concentração acima dele.
As distribuições por Grande Região também apresentam assimetria negativa, maior em valor absoluto nas regiões Norte e Nordeste, 1,17 e 0,84, respectivamente.
Os Gráficos 3.2, 3.3 e 3.4 apresentam informações referentes à média da nota final dos Participantes, desagregando os resultados de acordo com, respectivamente, as Grandes Regiões do país, a Categoria Administrativa e a Organização Acadêmica.
Os gráficos apresentam o valor da média das notas como uma barra e os extremos do intervalo de confiança de 95 \% como linhas verticais unidas por uma linha horizontal na forma da letra H maiúscula.
Considerando -se o gráfico de notas segundo Grande Região ( Gráfico 3.2 ), observa -se que existe diferença estatisticamente significativa, ao nível de 95 \%, entre a maior média, obtida na região Nordeste ( 57,8 ) e a menor, obtida na região Sudeste ( 48,7 ).
Levando -se em conta os estudantes por Categorias Administrativas ( Gráfico 3.3 ), observa -se que existe diferença estatisticamente significativa entre as médias das notas das IES Públicas e Privadas.
Pode ser observado ainda que a média dos alunos de IES Públicas ( 53,6 ) é maior que a dos alunos de IES Privadas ( 49,2 ).
No entanto, a diferença entre IES Públicas e Privadas é menor do que a diferença entre as regiões Nordeste e Sudeste ( 9,1 ).
Tendo como foco o Gráfico 3.4, que apresenta as notas médias das provas segundo a Organização Acadêmica, contata -se que as diferenças nas médias das notas dos estudantes provenientes de Universidades e Faculdades é estatisticamente significativa.
No entanto, as diferenças entre os Centros Universitários e as outras duas categorias não são estatisticamente significativas.
A maior média foi obtida pelos estudantes de Universidades, e a menor, pelos de Faculdades. 3.1.2 Estatísticas Básicas no Componente de Formação Geral A Tabela 3.2 apresenta as estatísticas básicas em relação ao componente da prova que avalia a Formação Geral dos estudantes concluintes.
Os alunos de todo Brasil obtiveram desempenho médio de 51,5 e o desvio padrão das notas dos estudantes foi 18,1.
A maior média foi obtida na região Nordeste ( 58,7 ), e a menor, na região Sudeste ( 49,5 ).
As demais médias foram : 51,7 na região Sul ; 54,2 na região Centro-Oeste ; e 55,0 na região Norte.
Já o maior desvio padrão foi obtido na região Sudeste ( 19,4 ) e o menor na região Centro-Oeste ( 15,3 ).
Os demais desvios padrões foram : 16,9 na região Norte ; 16,7 na região Nordeste e 17,0 na região Sul.
A maior nota no componente de Formação Geral da prova do ENADE foi obtida por pelo menos um aluno da região Sudeste ( 94,0 ) enquanto a menor nota máxima ocorreu nas regiões Norte e Nordeste ( 88,5 ).
Nas outras regiões, as notas máximas foram : 90,5 na região Sul e 93,0 na região Centro-Oeste.
A mediana do Brasil como um todo foi 53,5, sendo a menor mediana encontrada na O Gráfico 3.5 propicia a avaliação do desempenho dos estudantes no componente de Formação Geral a partir do histograma da distribuição das notas correspondentes.
A distribuição é unimodal, com moda em ( 50 ; 60 ], o mesmo intervalo modal encontrado para a prova como um todo.
Nota -se, ainda, que o Gráfico 3.1 ( distribuição das notas da prova ) é parecido com o agora analisado ( Gráfico 3.5 ), apesar de haver uma dispersão um pouco maior das notas do Componente de Formação Geral.
Isso se deve à diferença entre os desvios padrões : 15,7 para a nota da prova como um todo e 18,1 para o componente de Formação Geral.
Para o componente de Formação Geral, o coeficiente de assimetria da distribuição das notas dos estudantes, como na prova como um todo, também é negativo ( 0,52 ).
Por ser um coeficiente pequeno, a distribuição apresenta uma leve concentração à direita e cauda maior à esquerda.
Em todas as Grandes Regiões os histogramas também possuem assimetria negativa com coeficientes variando entre 0,32 e 0,71.
Nos Gráficos 3.6, 3.7 e 3.8 são apresentadas as informações referentes ao desempenho dos Concluintes no componente de Formação Geral, em diferentes agregações : Grande Região do país, Categoria Administrativa e Organização Acadêmica.
Observa -se pelo Gráfico 3.6 que existe diferença estatisticamente significativa entre as médias das notas das regiões Sudeste e Nordeste, menor e maior média, no Componente de Formação Geral, segundo a Grande Região do país.
Vemos que o intervalo de confiança é mais largo nas regiões Norte e Nordeste.
Este fato está relacionado com o fato do tamanho da população envolvida ser menor nessas regiões ( 116 presentes na região Norte e 158 na Nordeste ).
No Gráfico 3.7, que representa as notas médias no Componente de Formação Geral segundo Categoria Administrativa do país, observa -se que a diferença entre as médias não é significativa, apesar de os concluintes das IES Públicas terem obtido uma média maior do que os das IES Privadas.
Considerando -se o tipo de Organização Acadêmica, nota -se, no Gráfico 3.8, que a diferença entre as médias, mais uma vez, não é estatisticamente significativa.
Observa -se, ainda, que não existem diferenças estatisticamente significativas entre as m ; edia de alunos de Universidades, Centro Universitário e faculdades. 3.1.3 Estatísticas Básicas do Componente de Conhecimento Específico A Tabela 3.3 apresenta as estatísticas básicas referentes ao componente de Conhecimento Específico da área de Licenciatura em Artes Visuais.
A média do desempenho dos estudantes do Brasil como um todo foi 51,0.
A maior média foi obtida na região Nordeste ( 57,5 ), e a menor, na região Sudeste ( 48,5 ).
As demais médias foram : 54,3 na região Norte, 53,3 na região Centro-Oeste e 52,5 na região Sul.
Quanto à variabilidade das notas, o desvio padrão do Brasil como um todo foi 16,9, sendo o maior desvio padrão observado na região Sudeste ( 18,3 ) e o menor na região Centro-Oeste ( 14,6 ).
Os demais desvios foram : 15,1 das regiões Norte e Sul e 14,8 da região Nordeste.
A mediana das notas dos estudantes de todo o Brasil foi 52,8.
A maior mediana ocorreu na região Nordeste ( 59,5 ) e a menor na região Sudeste ( 50,9 ).
As demais medianas foram : 56,9 na região Norte ; 54,3 na região Sul ; e 53,5 na região Centro-Oeste.
A nota máxima do Brasil como um todo foi 90,9, sendo obtida por pelo menos um aluno da região Nordeste.
As demais notas máximas foram : 89,3 na região Sudeste ; 85,4 na região Centro-Oeste ; 84,9 na região Sul ; e 83,7 na região Norte.
A nota mínima, foi diferente de zero apenas na região Centro-Oeste. ( 11,1 ).
Assim como os Gráficos 3.1 e 3.5, o Gráfico 3.9, apresentado a seguir, proporciona uma visão da avaliação do desempenho de concluintes em relação ao componente de Conhecimento Específico, com um histograma da distribuição das notas correspondentes.
Como nos histogramas anteriores, esta é uma distribuição unimodal com intervalo modal em ( 50 ; 60 ].
O coeficiente de assimetria da distribuição das notas do componente de Conhecimento Específico é – 0,58.
Este coeficiente confirma, como mostra o histograma ( Gráfico 3.9 ), que a distribuição não é simétrica.
Ela possui uma cauda mais longa abaixo do intervalo modal e concentração a partir dele.
Na região Norte o coeficiente de assimetria é bem próximo de – 1, evidenciando que a distribuição das notas dos estudantes dessa região, no componente de Conhecimento Específico, é ainda mais assimétrica.
Oeste ; 84,9 na região Sul ; e 83,7 na região Norte.
A nota mínima, foi diferente de zero apenas na região Centro-Oeste. ( 11,1 ).
Assim como os Gráficos 3.1 e 3.5, o Gráfico 3.9, apresentado a seguir, proporciona uma visão da avaliação do desempenho de concluintes em relação ao componente de Conhecimento Específico, com um histograma da distribuição das notas correspondentes.
Como nos histogramas anteriores, esta é uma distribuição unimodal com intervalo modal em ( 50 ; 60 ].
O coeficiente de assimetria da distribuição das notas do componente de Conhecimento Específico é – 0,58.
Este coeficiente confirma, como mostra o histograma ( Gráfico 3.9 ), que a distribuição não é simétrica.
Ela possui uma cauda mais longa abaixo do intervalo modal e concentração a partir dele.
Na região Norte o coeficiente de assimetria é bem próximo de – 1, evidenciando que a distribuição das notas dos estudantes dessa região, no componente de Conhecimento Específico, é ainda mais assimétrica.
Os Gráficos 3.10, 3.11 e 3.12 apresentam uma comparação dos resultados em relação à Grande Região do país, à Categoria Administrativa e à Organização Acadêmica, agora levando em conta o desempenho de estudantes no componente de Conhecimento Específico da prova.
Pelo Gráfico 3.10, observa -se que existe diferença estatisticamente significativa entre as médias das notas no componente de Conhecimento Específico, da região Nordeste, onde a média foi a mais elevada, em relação às regiões Sudeste, Sul e Centro-Oeste.
A região Sudeste foi onde os concluintes obtiveram a menor média, sendo também a que possui intervalo de confiança mais estreito o que está relacionado ao fato de ter sido a região com o maior número de concluintes presentes ao exame ( 1.453 ).
Quanto à Categoria Administrativa ( gráfico 3.11 ), observa -se que, diferentemente do Componente de Formação Geral, existe diferença estatisticamente significativa entre as médias das IES Públicas e Privadas para o Componente de Conhecimento Específico.
No entanto, mais uma vez, a maior média foi obtida por alunos de IES Públicas de ensino.
Quanto ao Gráfico 3.12, observa -se que existe diferença estatisticamente significativa ao nível de 95 \% entre as notas no componente de Conhecimento Específico dos alunos de Universidade em relação aos outros tipos de Organização Acadêmica.
Sendo a média dos Concluintes das Universidades ( 52,0 ) maior do que a dos alunos de Centros Universitários ( 49,3 ) e de Faculdades ( 48,4 ).
As diferenças entre a média de Universidades e as duas outras categorias é estatisticamente significativa a 5 \%, mas a diferença entre Centros Universitários e Faculdades, não. 3.2 ANÁLISE DAS QUESTÕES OBJETIVAS 3.2.1 Componente de Formação Geral A Tabela 3.4 apresenta as estatísticas básicas relativas às oito questões objetivas do componente da prova que abrange a Formação Geral dos estudantes.
A média do Brasil foi 50,5.
A menor média foi encontrada na região Sudeste ( 49,1 ) e a maior na região Nordeste ( 58,6 ).
As demais médias foram : 57,0 na região Norte ; 51,3 na região Centro-Oeste ; e 50,1 na região Sul.
O desvio padrão do Brasil foi 19,7, sendo o maior desvio padrão encontrado na região Sudeste ( 20,7 ) e o menor na região Nordeste ( 16,7 ).
Os demais desvios foram : 17,7 na região Norte ; 18,4 na região Centro-Oeste ; e 18,9 na região Sul.
As medianas corresponderam a 50,0 para todo o Brasil e para as regiões Sudeste, Sul e Centro-Oeste.
Já para as regiões Norte e Nordeste a mediana foi mais elevada, 62,5.
A nota máxima, 100,0, foi obtida por pelo menos um aluno das regiões Nordeste, Sudeste e Centro-Oeste e foi igual a 87,5 nas regiões Norte e Sul.
Já a nota mínima foi 0,0 em todas as regiões.
A Tabela 3.5 apresenta os índices de facilidade e discriminação ( ponto bisserial ) das questões objetivas do componente de Formação Geral.
Quanto ao índice de facilidade, foram usadas as seguintes cores para diferenciar o nível de dificuldade da questão : • Azul para as questões classificadas com índice muito fácil ( índice \&gt ; =0,86 ), verde para as questões classificadas com índice fácil ( 0,61 a 0,85 ), amarelo para as questões classificadas com médio ( 0,41 a 0,60 ), vermelho para as questões classificadas com difícil ( 0,16 a 0,40 ) e roxo para as questões classificadas com muito difícil ( \&lt ; =0,15 ).
Já quanto ao índice de discriminação, foram usadas as seguintes cores para qualificar a questão : • As questões classificadas com índice fraco receberam a cor vermelho ( índice \&lt ; =0,19 ), as classificadas com médio receberam a cor amarelo ( 0,20 a 0,29 ), as classificadas com bom receberam a cor verde ( 0,30 a 0,39 ) e as classificadas com muito bom ( \&gt ; =0,40 ) receberam a cor azul.
As questões objetivas do componente de Formação Geral, segundo o índice de facilidade, foram assim avaliadas : das oito questões, nenhuma teve o índice de facilidade classificado como muito fácil.
Três questões foram tidas como fáceis, por terem índice de acertos situado na faixa entre 0,61 e 0,85 ( de 61,0 \% a 85,0 \% de acertos ).
Duas questões foram consideradas de dificuldade média, situando -se no intervalo entre 0,41 e 0,60 do índice de facilidade, ou seja, houve entre 41,0 \% e 60,0 \% de acertos.
Outras três questões foram classificadas na categoria difícil, situando -se no intervalo entre 0,16 e 0,40.
Por fim, nenhuma das questões apresentou menos de 15 \% de acertos, razão pela qual seriam classificadas como muito difíceis.
A Tabela 3.5 apresenta os índices de facilidade e discriminação ( ponto bisserial ) das questões objetivas do componente de Formação Geral.
Quanto ao índice de facilidade, foram usadas as seguintes cores para diferenciar o nível de dificuldade da questão :  Azul para as questões classificadas com índice muito fácil ( índice \&gt ; =0,86 ), verde para as questões classificadas com índice fácil ( 0,61 a 0,85 ), amarelo para as questões classificadas com médio ( 0,41 a 0,60 ), vermelho para as questões classificadas com difícil ( 0,16 a 0,40 ) e roxo para as questões classificadas com muito difícil ( \&lt ; =0,15 ).
Já quanto ao índice de discriminação, foram usadas as seguintes cores para qualificar a questão :  As questões classificadas com índice fraco receberam a cor vermelho ( índice \&lt ; =0,19 ), as classificadas com médio receberam a cor amarelo ( 0,20 a 0,29 ), as classificadas com bom receberam a cor verde ( 0,30 a 0,39 ) e as classificadas com muito bom ( \&gt ; =0,40 ) receberam a cor azul.
As questões objetivas do componente de Formação Geral, segundo o índice de facilidade, foram assim avaliadas : das oito questões, nenhuma teve o índice de facilidade classificado como muito fácil.
Três questões foram tidas como fáceis, por terem índice de acertos situado na faixa entre 0,61 e 0,85 ( de 61,0 \% a 85,0 \% de acertos ).
Duas questões foram consideradas de dificuldade média, situando -se no intervalo entre 0,41 e 0,60 do índice de facilidade, ou seja, houve entre 41,0 \% e 60,0 \% de acertos.
Outras três questões foram classificadas na categoria difícil, situando -se no intervalo entre 0,16 e 0,40.
Por fim, nenhuma das questões apresentou menos de 15 \% de acertos, razão pela qual seriam classificadas como muito difíceis.
Como já comentado, para análise das questões objetivas relativas à Formação Geral segundo o poder de discriminação, utilizou -se, o índice de discriminação ponto bisserial.
Nesta análise as questões foram assim avaliadas : cinco das oito questões apresentaram índices acima ou igual a 0,40 e, assim, foram classificadas com índice muito bom para esse grupo de estudantes.
Duas questões tiveram bom índice de discriminação, com valor entre 0,30 e 0,39 e uma outra apresentou um valor entre 0,20 e 0,29, sendo classificada com índice médio para esse grupo de estudantes.
Nenhuma questão teve nível fraco de discriminação.
O índice de facilidade variou de 0,18 a 0,74, e o de discriminação, de 0,24 a 0,54.
As cinco questões com índices de discriminação muito bom, as de números 1, 2, 3, 5 e 6, figuraram entre os diversos níveis de dificuldade desse conjunto : três classificadas na categoria fácil ( questões 1, 3 e 6 ) do índice de facilidade, uma na categoria médio ( questão 5 ) e uma na categoria difícil ( questão 2 ).
Em particular, a questão 3 foi a que apresentou o maior poder discriminatório, com índice 0,54, e foi a segunda mais fácil, com uma proporção de 0,69 acertos.
O máximo de acertos foi alcançado pela questão 6 com um índice de facilidade de 0,74.
A questão de número 4 apresentou um índice de facilidade de 0,56, ou seja, um pouco mais da metade dos estudantes conseguiu resolvê -la, dentro do universo de participantes.
Seu índice de discriminação foi bom ( 0,35 ).
Já a questão 8 obteve índice de discriminação médio, 0,24, e seu índice de facilidade foi difícil ( 0,18 ).
O Gráfico 3.13, para exemplificar, analisa o comportamento da questão de número 3 de Formação Geral.
Trata -se da segunda questão mais fácil e a que obteve o maior índice de discriminação dessa parte da prova.
Neste gráfico, cada uma das cinco curvas representa o percentual de respostas em determinada alternativa da questão, em função do número de acertos dos estudantes nesta parte da prova ( Formação Geral/Múltipla Escolha ), antes de possíveis eliminações pelo critério do ponto bisserial.
Em princípio, a soma das escolhas possíveis deveria ser igual a 100 \%.
Não é este o caso, notadamente para o número zero de acertos, pois não aparecem no gráfico as questões deixadas em branco ou com múltiplas respostas.
Como foram oito as questões, os valores variam de 0 a 8 acertos.
A curva em vermelho corresponde à alternativa E, a correta para esta questão.
Assim, observa -se que entre os estudantes com menor número de acertos, nessa parte do exame, a situação mais frequente foi a escolha de uma das alternativas incorretas : a alternativa A ( em azul ) ou D ( em roxo ).
Na medida em que o número de acertos aumenta, indicando desempenho melhor nesta parte da prova, aumenta concomitantemente a proporção de estudantes que selecionaram a alternativa correta E, atingindo 100 \% para os estudantes com 7 acertos.
Essa análise permite verificar como a questão discriminou os grupos de desempenho, justificando o alto índice obtido na questão.
Os gráficos relativos às demais questões de Formação Geral constam do Anexo I. 3.2.2 Componente de Conhecimento Específico A Tabela 3.6 apresenta as estatísticas básicas em relação às questões objetivas do componente de Conhecimento Específico da prova, por Grande Região.
A média do Brasil deste componente foi de 54,6.
A menor média foi observada na região Sudeste ( 52,0 ) e a maior na região Nordeste ( 61,8 ).
O desvio padrão de todo o Brasil foi 18,2, sendo o menor desvio padrão encontrado na região Centro-Oeste ( 15,8 ) e o maior na região Sudeste ( 19,6 ).
A mediana de todo o Brasil foi 56,5.
As medianas por região foram : 60,9 na região Norte, 65,2 na região Nordeste e 56,5 nas outras três grandes regiões.
A nota máxima da prova ( 95,7 ) foi obtida no componente de Conhecimento específico por pelo menos um aluno das regiões Nordeste e Sudeste.
A nota mínima só foi diferente de zero ( 13,0 ) na região Centro-Oeste.
A Tabela 3.7 apresenta os índices de facilidade e discriminação ( ponto bisserial ) das questões objetivas do componente de Conhecimento Específico.
Para facilitar a diferenciação das questões usaram -se as mesmas cores da Tabela 3.5 para as diferentes classificações dos índices de facilidade e de discriminação.
Dentre as questões objetivas da parte da prova relativa ao componente de Conhecimento Específico, nenhuma questão foi anulada pela Comissão.
Desse modo, a classificação quanto ao índice de facilidade foi estabelecida com base nas 27 questões.
A partir dos índices obtidos, pode -se concluir que as questões objetivas da prova estavam distribuídas em torno de uma dificuldade média com dez questões nesta categoria.
As outras 17 questões estavam quase que igualmente distribuídas acima e abaixo deste valor : nove classificadas como difícil e oito como fácil.
Não houve questão classificada como muito difícil, nem tampouco muito fácil. nove classificadas como difícil e oito como fácil.
Não houve questão classificada como muito difícil, nem tampouco muito fácil.
Já quanto aos índices de discriminação das questões objetivas do componente de Conhecimento Específico da prova, tem -se como resultado a seguinte classificação : oito das 27 questões foram consideradas como boas, enquanto dez delas tiveram índice de discriminação muito bom.
Dentre as demais, cinco delas foram classificadas como médias e quatro como fracas, sendo nove, por conseguinte, a quantidade de questões nos dois patamares mais baixos de discriminação.
Constata -se, assim, que a prova – no que se refere ao componente de Conhecimento Específico – possuía boa capacidade de discriminar entre aqueles que dominam ou não o conteúdo.
Dentre as dez questões que alcançaram os maiores índices de discriminação, as de números 10, 14, 15, 25, 29 e 33 foram classificadas como fáceis e as demais ( 19, 20, 22 e 23 ) com grau de dificuldade médio.
A questão de número 11 foi a mais difícil dentre as 27 questões específicas, classificada como difícil pelo índice de facilidade com apenas 19,0 \% de acertos.
Essa questão apresentou poder discriminatório médio, 0,25.
Quatro questões, as de números 9, 26, 34 e 35 apresentaram baixo poder discriminatório, classificadas como fracas pelo índice de discriminação e foram eliminadas do cômputo da nota final.
Estas questões foram todas classificadas como difíceis com respeito ao índice de facilidade.
A título de exemplo das análises do comportamento das questões objetivas, o Gráfico 3.14 analisa a questão 33 do componente de Conhecimento Específico.
Esta foi a quarta questão mais fácil da prova, apresentando índice de facilidade 0,80, ou seja, 80,0 \% dos estudantes assinalaram acertadamente a opção D, correspondente ao gabarito.
Seu índice de discriminação foi igual a 0,53, classificado como muito bom.
Neste gráfico, cada uma das cinco curvas representa o percentual de respostas em determinada alternativa da questão 33 em função do número de acertos dos estudantes nesta parte da prova, antes de possíveis eliminações de questões pelo critério do ponto bisserial.
A alternativa correta D, representada no gráfico pela curva em roxo, foi escolhida em maiores proporções pelos alunos com desempenho melhor nesta parte da prova.
Já as alternativas incorretas, também denominadas distratores, foram selecionadas principalmente por aqueles com notas mais baixas.
Em princípio, a soma das escolhas possíveis deveria ser igual a 100 \%.
Não é este o caso, notadamente para o número zero de acertos, pois aqueles com nota zero, na sua maioria deixaram esta questão em branco ou marcaram mais de uma alternativa, comportamento considerado inválido.
A proporção de alunos que selecionou a resposta correta D aumenta gradativamente, chegando a atingir 100 \% para mais de 21 acertos, enquanto a proporção dos que escolheram alternativas incorretas, principalmente as alternativas E – curva em vermelho e A, curva em azul, decai, como função do número de acertos nesta parte da prova.
Os gráficos relativos às demais questões do componente de Conhecimento Específico constam do Anexo I. 3.3 ANÁLISE DAS QUESTÕES DISCURSIVAS 3.3.1 Componente de Formação Geral As análises dos resultados de desempenho dos estudantes de Licenciatura em Artes Visuais nas duas questões discursivas relativas à Formação Geral encontram -se na Tabela 3.8 e no Gráfico 3.15.
Na tabela 3.8 observa -se que a nota média, considerando -se os concluintes de todo o Brasil ( 53,0 ) nas questões discursivas do Componente de Formação Geral, foi mais alta do que a média obtida nas questões objetivas ( 50,5 ) do mesmo Componente.
A mediana confirma o melhor desempenho dos alunos de todo o Brasil nas questões discursivas do componente de Formação Geral.
Enquanto essa estatística foi de 55,0 para questões discursivas, para as questões objetivas a mediana foi 50,0.
Pode -se notar um aumento do desvio padrão de 19,7 nas questões objetivas do componente de Formação Geral dos alunos de todo o Brasil, para 26,8 nas questões discursivas do mesmo componente.
A mediana de todo o Brasil, neste componente, foi 55,0, sendo a maior mediana encontrada na região Nordeste ( 65,0 ) e a menor nas regiões Norte e Sudeste ( 50,0 ).
A nota máxima ( 100,0 ) e a nota mínima ( 0,0 ) foram iguais em todas as regiões do Brasil.
O Gráfico 3.15 representa a distribuição das notas nas questões discursivas no componente de Formação Geral.
A moda dessa distribuição ocorre no intervalo ( 40 ; 50 ], mas observam -se outras modas locais nos intervalos [ 0 ; 10 ] e ( 70 ; 80 ].
O coeficiente de assimetria é negativo ( – 0,46 ) e o mesmo ocorre para todas as regiões, sendo a assimetria mais acentuada na região Nordeste ( – 0,76 ).
Na sequência, os resultados verificados para cada uma das questões discursivas de Formação Geral serão apresentados, estabelecendo -se relações com os conteúdos abordados em cada uma delas.
Os comentários da Banca de docentes corretores a respeito do observado na correção das respostas dos estudantes, suas impressões e conclusões, serão apresentados junto à análise de cada questão.
Cumpre esclarecer que, tendo em vista que as questões discursivas de Formação Geral são padronizadas, ou seja, constam de todas as provas, os comentários da Banca são os mesmos para todas as carreiras acadêmicas, sendo direcionados a todos os estudantes que participaram do ENADE/2011.
A seguir, serão analisados os desempenhos da Área de Licenciatura em Artes Visuais nas duas questões discursivas de Formação Geral do ENADE/2011, comparando os resultados obtidos com comentários para cada questão. 3.3.1.1 Análise da Questão Discursiva 1 do Componente de Formação Geral Os dados de Licenciatura em Artes Visuais, obtidos a partir das respostas à questão 1, encontram -se na Tabela 3.9 e no Gráfico 3.16.
Nessa questão – de melhor desempenho dentre as duas de Formação Geral – os alunos de todo Brasil tiveram média, 58,3.
A maior média para a questão 1 foi obtida na região Nordeste ( 68,0 ), e a menor, na região Sudeste ( 55,1 ).
Quanto à variabilidade das notas, o desvio padrão de todo o Brasil foi 34,7.
O menor desvio padrão foi obtido na região Centro-Oeste ( 31,0 ) e o maior desvio padrão foi obtido na região Sudeste ( 35,7 ).
As medianas e as notas máximas da questão discursiva 1 foram : 60,0 nas regiões Norte e Sudeste ; 65,0 na região Sul ; e 70,0 nas regiões Nordeste e Centro-Oeste.
A nota mínima ( 0,0 ) e a nota máxima ( 100,0 ) foram obtidas em todas as regiões do país, sem exceção.
O Gráfico 3.16 mostra a distribuição das notas na questão discursiva 1 do componente de Formação Geral.
Observa -se que a maior frequência está no intervalo ( 90 ; 100 ], o que confirma um bom desempenho dos alunos nesta questão. 3.3.1.2 Comentários sobre a correção das respostas à Questão Discursiva 1 De maneira geral, a aplicação da chave de correção da questão 1 de Formação Geral não apresentou qualquer dificuldade digna de menção.
Para isso contribuíram, indubitavelmente, os ajustes feitos após a correção da amostra e a reunião entre todos os membros da Banca de docentes corretores.
As poucas dúvidas, todas pontuais, apresentadas pelos corretores, foram acompanhadas e respondidas pela coordenação e subcoordenação da correção das questões de Formação Geral, por meio da ferramenta de Gerenciamento de Dúvidas do Sistema de Correção On-line.
Não houve registro de qualquer ocorrência que pusesse em xeque o padrão de resposta ou a efetividade e a adequação da chave de correção.
Explica -se : trata -se de questão com comando claro, direto e objetivo ( solicitava -se, basicamente, três vantagens justificadas de cursos a distância ), cujas respostas foram corrigidas por meio da aplicação de um chave de correção testada e aprovada previamente.
Havia absoluta clareza quanto aos critérios de avaliação da correspondência entre as respostas dos estudantes e as possibilidades de vantagens de cursos a distância admitidas como corretas no padrão de resposta oficial, além de gradações explícitas ( e fáceis de aplicar ) dos diferentes níveis de pontuação previstos.
Felizmente, portanto, não há reparo a registrar em relação à facilidade de aplicação do padrão de resposta e da chave de correção, e nem em relação à atribuição dos diferentes níveis de pontuação previstos.
Todas as dificuldades que poderiam ter obstado a correta aplicação do padrão de resposta oficial e da respectiva chave de correção foram evitadas por meio dos ajustes feitos após a correção da amostra e fartamente debatidos com toda a Banca.
Digno de nota é que quantidade tão significativa de profissionais envolvidos na mesma tarefa – tanto para a questão 1 quanto para a 2 – tenha apresentado tão poucas dificuldades na execução da correção, em termos proporcionais.
Em suma, a correção da questão 1 da prova de Formação Geral do ENADE 2011 foi exemplarmente bem planejada, servindo -se de padrão de resposta muito bem adequado à questão proposta.
Quanto ao tema desta questão, em particular, a Banca verificou que uma parcela significativa de estudantes evocou experiências bastante concretas e próximas de sua realidade.
Houve várias respostas que indicavam uma vivência pessoal de ensino superior na modalidade Educação a Distância ( EaD ), evidenciando o tom de depoentes nos textos apresentados.
Os estudantes, em sua grande maioria, utilizaram parte considerável do espaço de 15 linhas disponíveis para a resposta – e outra parcela menos significativa dos que não o fizeram demonstraram notável capacidade de atender ao comando da questão de maneira objetiva, curta e, via de regra, correta.
Registre -se, ainda, que foram relativamente poucos os casos de respostas que tenham passado ao largo do tema em pauta na questão.
As capacidades de leitura, de compreensão do comando proposto e de expressão escrita que os estudantes avaliados na edição 2011 do ENADE foram satisfatórias.
Não obstante, seria leviano perder de vista que a qualidade dos textos redigidos em resposta às questões discursivas do Exame ainda está muito aquém do que se espera de concluintes de cursos de ensino superior de todas as regiões do país.
Quanto ao conteúdo das respostas, a Banca constatou boa capacidade, por parte da maioria dos estudantes, de compreensão do tema e do comando da questão.
Foram relativamente poucos os casos de respostas que deixaram de enumerar vantagens da modalidade EaD, e proporcionalmente escassos os estudantes que citaram vantagens não previstas no padrão de resposta.
Foi frequente, a tentativa direta de atender ao comando da questão.
Os erros mais comuns, em relação ao padrão de respostas e à grade de correção, foram fruto do desdobramento em vários “ itens ” daquilo que, de acordo com o padrão de respostas oficial, representava uma única vantagem.
Destacaram -se, neste caso, as respostas que apontavam a flexibilidade de horário e/ou local como duas vantagens distintas daquela modalidade de ensino.
Quanto aos diferentes níveis de pontuação previstos, a maior causa de baixas pontuações foi a ausência de justificativas, e mesmo de argumentação, para uma ou mais das vantagens enumeradas.
Isso demonstra não apenas a objetividade das respostas, coerente com a objetividade do comando da questão ( “ enumere três vantagens de um curso a distância ” ), mas também certa dificuldade de formulação plena de um texto, ou ao menos de parágrafos, em formato dissertativo – mesmo diante de uma média de 5 linhas disponíveis para cada vantagem a enumerar.
Os acertos mais comuns, ou seja, os “ itens ” do padrão de respostas mais frequentemente mencionados foram : ( 1 ) a flexibilidade de horário e/ou local ; ( 2 ) a capilaridade do ensino a distância ; ( 3 ) a democratização do acesso à educação de qualidade ; e ( 4 ) os custos menores que os de cursos presenciais.
Dentre as vantagens previstas no padrão de respostas que foram menos citadas, destacam -se a inclusão de pessoas com comprometimento motor, a qualificação de professores e a troca de experiências entre os participantes.
Foram muito comuns, no entanto, as menções à supostamente maior facilidade de acesso a professores ou/e tutores em cursos superiores a distância. 3.3.1.3 Análise da Questão Discursiva 2 do Componente de Formação Geral A Tabela 3.10 mostra que o desempenho dos estudantes na questão 2 ( média 47,2 ) foi inferior ao obtido na questão de número 1 ( média 58,3 ).
A região Centro-Oeste foi aquela onde a média, nessa questão, foi maior ( 53,4 ), e a de menor média foi a região Sudeste ( 44,3 ).
Quanto à variabilidade das notas, o desvio padrão de todo o Brasil foi 29,9, inferior ao obtido na questão de número 1 ( 34,7 ), ou seja, além da média ser mais baixa houve menos variabilidade.
O maior desvio nessa questão foi obtido na região Norte ( 31,2 ), enquanto o menor foi obtido na região Centro-Oeste ( 27,1 ).
A mediana comprova a diferença de desempenho entre as duas questões discursivas do componente de Formação Geral.
As medianas foram, respectivamente, 60,0 e 50,0 nas questões 1 e 2.
A menor mediana ocorreu na região Sudeste ( 50,0 ), enquanto a maior ( 60,0 ) ocorreu nas regiões Norte, Nordeste e Centro-Oeste.
A nota máxima ( 100,0 ) foi obtida em todas as regiões do Brasil e a nota mínima ( 0,0 ) também. respostas oficial, representava uma única vantagem.
Destacaram -se, neste caso, as respostas que apontavam a flexibilidade de horário e/ou local como duas vantagens distintas daquela modalidade de ensino.
Quanto aos diferentes níveis de pontuação previstos, a maior causa de baixas pontuações foi a ausência de justificativas, e mesmo de argumentação, para uma ou mais das vantagens enumeradas.
Isso demonstra não apenas a objetividade das respostas, coerente com a objetividade do comando da questão ( “ enumere três vantagens de um curso a distância ” ), mas também certa dificuldade de formulação plena de um texto, ou ao menos de parágrafos, em formato dissertativo – mesmo diante de uma média de 5 linhas disponíveis para cada vantagem a enumerar.
Os acertos mais comuns, ou seja, os “ itens ” do padrão de respostas mais frequentemente mencionados foram : ( 1 ) a flexibilidade de horário e/ou local ; ( 2 ) a capilaridade do ensino a distância ; ( 3 ) a democratização do acesso à educação de qualidade ; e ( 4 ) os custos menores que os de cursos presenciais.
Dentre as vantagens previstas no padrão de respostas que foram menos citadas, destacam -se a inclusão de pessoas com comprometimento motor, a qualificação de professores e a troca de experiências entre os participantes.
Foram muito comuns, no entanto, as menções à supostamente maior facilidade de acesso a professores ou/e tutores em cursos superiores a distância. 3.3.1.3 Análise da Questão Discursiva 2 do Componente de Formação Geral A Tabela 3.10 mostra que o desempenho dos estudantes na questão 2 ( média 47,2 ) foi inferior ao obtido na questão de número 1 ( média 58,3 ).
A região Centro-Oeste foi aquela onde a média, nessa questão, foi maior ( 53,4 ), e a de menor média foi a região Sudeste ( 44,3 ).
Quanto à variabilidade das notas, o desvio padrão de todo o Brasil foi 29,9, inferior ao obtido na questão de número 1 ( 34,7 ), ou seja, além da média ser mais baixa houve menos variabilidade.
O maior desvio nessa questão foi obtido na região Norte ( 31,2 ), enquanto o menor foi obtido na região Centro-Oeste ( 27,1 ).
A mediana comprova a diferença de desempenho entre as duas questões discursivas do componente de Formação Geral.
As medianas foram, respectivamente, 60,0 e 50,0 nas questões 1 e 2.
A menor mediana ocorreu na região Sudeste ( 50,0 ), enquanto a maior ( 60,0 ) ocorreu nas regiões Norte, Nordeste e Centro-Oeste.
A nota máxima ( 100,0 ) foi obtida em todas as regiões do Brasil e a nota mínima ( 0,0 ) também.
O Gráfico 3.17 mostra a distribuição das notas na questão discursiva 2 do componente de Formação Geral.
Este gráfico evidencia que mais do que 15 \% dos alunos deixaram a questão 2 em branco, e quase 5 \% deles receberam nota zero.
Desconsiderando estes casos, o histograma apresenta moda no intervalo ( 50 ; 60 ].
Nota -se, ainda, que a distribuição possui assimetria negativa ( coeficiente – 0,34 ).
Este tipo de assimetria também ocorre para a distribuição por região, já que os coeficientes de assimetria variam entre – 0,20 na região Sudeste e – 0,67 na região Centro-Oeste. 3.3.1.4 Comentários sobre a correção das respostas à Questão Discursiva 2 Diferentemente da questão 1, a questão 2, cujo tema aborda políticas públicas para a erradicação do analfabetismo e para o aumento da empregabilidade, permitiu mais amplitude na elaboração das respostas dos estudantes, uma vez que solicita a apresentação de uma proposta.
À falta de objetividade técnica na elaboração da questão – sobretudo no seu comando – correspondeu uma miríade de possibilidades interpretativas.
Registre -se que não raras foram as respostas que se utilizaram de lugares comuns e exortações religiosas/humanitárias/cívicas, contudo, na grande maioria dos casos, essas exortações foram usadas como um encerramento do texto e não comprometeram a resposta.
Muito frequente também foi a confusão entre política pública e política partidária, bem como alguns poucos, confundiam programas educacionais com programas educativos nas TVs.
Assim, destacaram -se como propostas/programas mais recorrentes : • prosseguimento das políticas já existentes, mas com o aumento dos investimentos, normalmente sugerindo a ampliação da bolsa família, da Educação de Jovens e Adultos ( EJA ) e da EaD ; • valorização do magistério, construção de escolas e melhoria das atuais ; • parceria do governo com empresas para manutenção de salas de aula visando a alfabetização de seus funcionários, oferecimento de estágios e redução da carga horária em troca de isenção de impostos ; • parcerias com igrejas e ONGs para criação de espaços de alfabetização ; • escolas itinerantes e alfabetizadores em domicílio, principalmente para pessoas com dificuldade de locomoção, como os idosos e deficientes físicos ; • erradicação do trabalho infantil ; • vinculação da bolsa família não apenas à frequência, mas também e, principalmente, ao resultado obtido pelo aluno na escola ; • revisão das políticas atuais, sendo a mais frequente a extinção da aprovação automática ; • críticas consistentes em relação a modelos didáticos considerados inadequados e desestimulantes para a educação não só de adultos, mas de pessoas de todas as idades.
Destacam -se, ainda, outras sugestões apresentadas : • Creches nas escolas onde os pais estão sendo alfabetizados.  Diminuição de duas horas na jornada de trabalho em empresas para funcionários não alfabetizados, para que possam frequentar a escola.  Campanhas educativas vinculadas aos meios de comunicação.  Aumento do número de escolas noturnas.
Formação específica para professores alfabetizadores.
Quanto à relação entre o analfabetismo e a empregabilidade, deve -se sublinhar que nem todos os estudantes estabeleceram claramente o vínculo entre essas duas situações sociais.
Alguns falaram separadamente de uma e de outra.
Mas a maioria fez referência à necessidade de estudo para “ conseguir um bom emprego com um bom salário ”.
Alguns, em menor número, estabeleceram de forma bastante interessante a questão histórica para a situação do Nordeste ; e, também, a relação entre a pessoa analfabeta/com pouca instrução e a desempregada/empregada em serviços mais pesados e pior remunerados, que não teriam condições de educar seus filhos que, por sua vez, também não teriam melhores oportunidades no mercado de trabalho, identificando a formação de um círculo vicioso e a necessidade de sua interrupção.
Muitos afirmaram que o analfabetismo não é o único responsável pelo desemprego e sim a má distribuição de renda.
Grande também foi o número de alunos que criticaram o resultado da pesquisa, afirmando que a mesma é enganosa, uma vez que considera alfabetizada a pessoa que “ desenha seu nome ”, sendo comum associarem a atual política de alfabetização com ganhos eleitorais.
Foi comum a resposta incompleta, em que o estudante fez apenas a análise das desigualdades/crítica do quadro apresentado, ou só apresentou proposta.
Alguns se limitaram a responsabilizar o governo referindo -se de forma bastante genérica à questão da “ educação ” e “ profissional ”.
Entre os equívocos que mais se repetiram, destacou -se a simples análise dos dados apresentados na tabela que consta do enunciado.
Alguns textos, inclusive, estavam corretos, sem que, no entanto, fosse respondida a pergunta.
Da mesma forma, quando os estudantes partiram para a segunda parte da questão, fizeram referência a vários programas já existentes ou simplesmente disseram que é muito importante que existam projetos para a educação.
Outro equívoco recorrente foi a análise da educação no Brasil como um todo.
Essa análise, apesar de correta, não se referia especificamente à questão do analfabetismo, o que, por conseguinte, levava a sugestões que não eram direcionadas à erradicação do analfabetismo, e sim à melhora da educação no Brasil.
Sendo assim, foram apresentadas propostas como o aumento do número de faculdades, o reforço de alunos do Ensino Médio, a criação de escolas técnicas, etc.
Por outro lado, houve um grande número de redações bem escritas e precisas, no que se refere ao que foi exigido pela questão.
Foram análises equilibradas e sensatas, correlacionando a problemática do analfabetismo com o desemprego e a apresentação de sugestões bastante consistentes. 3.3.2 Componente de Conhecimento Específico Na parte da prova relativa às questões discursivas no componente de Conhecimento Específico ( Tabela 3.11 ), constata -se que a média foi mais baixa do que para as questões discursivas do componente de Formação Geral.
Enquanto no componente de Formação Geral a média para estudantes de Licenciatura em Artes Visuais de todo o Brasil foi 53,0, na parte de Conhecimento Específico a média foi 30,3.
A maior média deste componente foi obtida pelos estudantes da região Centro-Oeste ( 33,0 ), e a menor, pelos da região Sudeste ( 28,3 ).
Quanto à variabilidade das notas, o desvio padrão de todo o Brasil foi 19,4.
O maior desvio padrão foi encontrado na região Norte ( 21,2 ), e o menor, na região Sul ( 18,4 ).
A maior nota máxima foi obtida na região Nordeste ( 88,3 ) enquanto a menor nota máxima foi encontrada na região Norte ( 75,0 ).
Além disso, a nota mínima ( 0,0 ) foi obtida por pelo menos um aluno em todas as regiões do Brasil, sem exceção.
A mediana do Brasil como um todo foi 31,7, bem mais baixa do que a mediana das notas do componente de Formação Geral, que foi 55,0.
Nas regiões Norte e Sudeste a mediana foi 30,0, na região Nordeste 32,5, na região Sul 33,3 e, por fim, no Centro-Oeste a mediana foi 35,0.
O Gráfico 3.18 representa a distribuição das notas nas questões discursivas no componente de Conhecimento Específico.
Essa distribuição apresenta duas modas, uma no intervalo [ 0 ; 10 ] e outra no intervalo ( 30 ; 40 ].
Nota -se que há um decrescimento rápido a partir do intervalo ( 30,40 ].
O coeficiente de assimetria do histograma das notas do Brasil como um todo é igual a zero, sendo que nas regiões Sul e Centro-Oeste é negativo ( – 0,08 e – 0,12, respectivamente ) e na demais regiões é positivo, mas também bem próximo de zero. 3.3.2.1 Análise da Questão Discursiva 3 do Componente de Conhecimento Específico Na questão 3, cujos resultados aferidos encontram -se descritos na Tabela 3.12, a média dos estudantes de todo o Brasil foi 33,0.
A menor média nessa questão foi obtida pelos alunos da região Norte ( 30,2 ), enquanto a maior média foi obtida na região Nordeste ( 37,0 ).
Quanto à variabilidade das notas, o desvio padrão de todo o Brasil foi 22,2.
O maior desvio padrão foi obtido na região Norte ( 24,0 ), enquanto o menor foi obtido na região Sul ( 21,4 ).
A nota máxima, 100,0 pontos, foi alcançada por pelo menos um aluno em cada uma das seguintes regiões : Sudeste e Nordeste.
As demais regiões tiveram como nota máxima 95,0.
A mediana do Brasil como um todo foi 40,0, sendo a mesma nas regiões Sul e Centro-Oeste, enquanto na região Nordeste foi mais alta, 45,0.
Nas demais regiões, a mediana foi 35,0.
A nota mínima ( 0,0 ) foi obtida em todas as regiões do Brasil, sem exceção.
O Gráfico 3.19 mostra a distribuição das notas na questão discursiva 3, do componente de Conhecimento Específico.
Essa distribuição é bimodal com moda nas questões em branco e no intervalo ( 40 ; 50 ] e o percentual de notas que ficaram nestes intervalos equivale a aproximadamente 50 \% do total de presentes.
As demais notas se distribuem com pequena assimetria negativa ( coeficiente igual a – 0,19 ) e o mesmo ocorre para as demais regiões ( coeficientes variando entre – 0,06 e – 0,35 ), a menos da região Norte para a qual o coeficiente de assimetria resultou positivo ( 0,09 ), ainda que pequeno. 3.3.2.2 Comentários sobre a correção das respostas à Questão Discursiva 3 O item aborda o conceito de artesania, em vez de artesanato.
A questão elaborada discute o fazer artístico a partir de uma análise comparativa dos procedimentos criativos de duas artistas plásticas, Beatriz Milhazes e Magadalena Atria.
A Banca considera que, independente de o estudante conhecer ou não as artistas citadas, o que mais importava na questão era o fato de as obras dessas artistas dialogarem com os próprios limites da pintura, na medida em que seus trabalhos apontam para modos de feitura, como a colagem, a \textbf{reprodução}, a incorporação do diálogo com as artes decorativas, que aprofundam o conceito de artesania, de modo a se contrapor a uma outra tendência da arte contemporânea, marcada pelo deslocamento contextual do objeto, enquanto forma de criação poética.
Deste modo, a questão solicitava ao estudante uma reflexão sobre estas características do fazer artístico contemporâneo.
As artistas citadas no enunciado da questão são exemplares, na medida em que trabalham com materiais e procedimentos que estão muito próximos dos materiais que o arte-educador utiliza comumente em suas atividades letivas e criativas ( massas para modelar, colagem, plástico, tecidos, entre outros ).
A partir dessa perspectiva, a questão potencializa a construção de uma ideia que não limita respostas a uma esfera estritamente teórica e conceitual.
Ao contrário, seu eixo de análise abre a possibilidade de um discurso que utilize elementos inseridos no movimento cotidiano de cada um, ou seja, em atividades vivenciadas em sala de aula e em ambientes onde se disseminam atividades artísticas, entre outros.
Nesse sentido, fica evidente que os participantes do ENADE pensaram suas respostas a partir de três quadrantes : o teórico, o conceitual e o empírico.
No caso do último quadrante, foi visível ter sido ele o eleito para uma solução discursiva num percentual significativo de respostas.
Apesar da reconsideração do Padrão de Resposta, tendo em vista contemplar uma maior diversidade de propostas, consideramos que a média dos estudantes avaliados permaneceu abaixo do desejado.
Nesta questão, predominaram as respostas classificadas como Insuficientes e Medianas.
As notas situadas entre o intervalo de 65 e 100 pontos constituíram uma parcela pouco significativa do total.
Conforme já havia ocorrido na Correção Amostral, por sua natureza e sua maior proximidade à experiência dos estudantes de Arte, a Questão 3 foi, dentre as três avaliadas, a que recebeu um menor percentual de respostas em branco.
No processo de correção, houve respostas consideradas “ ótimas ”, por terem correspondido ao Padrão de Resposta e por manifestar uma compreensão crítica acerca da questão, propondo, inclusive, uma mudança conceitual.
As respostas medianas foram aquelas em que se observou uma menor clareza do ponto de vista conceitual.
Mas, de um modo geral, notou -se que alguns conteúdos fundamentais para a formação do arte-educador foram assimilados pelos alunos, como é o caso do conceito de linguagem visual e expressividade.
Verificou -se, nesse tipo de resposta, uma crescente limitação da linguagem específica do estudante acerca dos conteúdos e conceitos.
De uma forma geral, as respostas que receberam baixa pontuação optaram por condensar os itens da questão num único discurso.
Em momentos isolados percebeu -se um aprofundamento discursivo, porém foi mais frequente a adoção de uma determinada linha de raciocínio que contemplou um item da questão em detrimento do outro.
Por conta dessa escolha, parece não ter havido assimilação clara, por parte dos estudantes, de como deveriam se constituir suas respostas, pois abordaram cada item de maneira rarefeita.
As respostas merecedoras da nota zero foram aquelas que, além de serem muito breves, não revelaram nenhuma formulação conceitual e tenderam a se desviar do tema da questão. modelar, colagem, plástico, tecidos, entre outros ).
A partir dessa perspectiva, a questão potencializa a construção de uma ideia que não limita respostas a uma esfera estritamente teórica e conceitual.
Ao contrário, seu eixo de análise abre a possibilidade de um discurso que utilize elementos inseridos no movimento cotidiano de cada um, ou seja, em atividades vivenciadas em sala de aula e em ambientes onde se disseminam atividades artísticas, entre outros.
Nesse sentido, fica evidente que os participantes do ENADE pensaram suas respostas a partir de três quadrantes : o teórico, o conceitual e o empírico.
No caso do último quadrante, foi visível ter sido ele o eleito para uma solução discursiva num percentual significativo de respostas.
Apesar da reconsideração do Padrão de Resposta, tendo em vista contemplar uma maior diversidade de propostas, consideramos que a média dos estudantes avaliados permaneceu abaixo do desejado.
Nesta questão, predominaram as respostas classificadas como Insuficientes e Medianas.
As notas situadas entre o intervalo de 65 e 100 pontos constituíram uma parcela pouco significativa do total.
Conforme já havia ocorrido na Correção Amostral, por sua natureza e sua maior proximidade à experiência dos estudantes de Arte, a Questão 3 foi, dentre as três avaliadas, a que recebeu um menor percentual de respostas em branco.
No processo de correção, houve respostas consideradas “ ótimas ”, por terem correspondido ao Padrão de Resposta e por manifestar uma compreensão crítica acerca da questão, propondo, inclusive, uma mudança conceitual.
As respostas medianas foram aquelas em que se observou uma menor clareza do ponto de vista conceitual.
Mas, de um modo geral, notou -se que alguns conteúdos fundamentais para a formação do arte-educador foram assimilados pelos alunos, como é o caso do conceito de linguagem visual e expressividade.
Verificou -se, nesse tipo de resposta, uma crescente limitação da linguagem específica do estudante acerca dos conteúdos e conceitos.
De uma forma geral, as respostas que receberam baixa pontuação optaram por condensar os itens da questão num único discurso.
Em momentos isolados percebeu -se um aprofundamento discursivo, porém foi mais frequente a adoção de uma determinada linha de raciocínio que contemplou um item da questão em detrimento do outro.
Por conta dessa escolha, parece não ter havido assimilação clara, por parte dos estudantes, de como deveriam se constituir suas respostas, pois abordaram cada item de maneira rarefeita.
As respostas merecedoras da nota zero foram aquelas que, além de serem muito breves, não revelaram nenhuma formulação conceitual e tenderam a se desviar do tema da questão.
Foi possível encontrar respostas nas quais ficou claro que o padrão proposto foi alcançado total ou parcialmente.
Quanto às parciais, um grande percentual de respostas foi recorrentemente simplista.
Houve um desdobramento marcante desse efeito, com empobrecimento das respostas e a criação de um lugar comum.
Com isso, não se pode afirmar que não tenham atingido ao objetivo estipulado pelo Padrão Resposta, mas apenas sinalizam : 1 ) a priorização dada a um dos itens da questão ; 2 ) o surgimento de um consenso de resposta representado por propostas de ensino de artes ora desconexas, ora estruturas com coerência, demonstrando preocupação apenas com as etapas do processo e com as formas de abordagem.
Além dessa observação, não se pode deixar de mencionar os conflitos no entendimento da resposta devido às limitações gramáticas e ortográficas de determinados estudantes. 3.3.2.3 Análise da Questão Discursiva 4 do Componente de Conhecimento Específico O Gráfico 3.20, representa a distribuição de notas na questão discursiva 4, no componente de Conhecimento Específico.
O percentual de presentes que deixou esta questão em branco foi alto, em torno de 32,5 \%, a moda da distribuição.
Dentre os que tentaram resolver a questão, pouco mais do que 10 \% teve nota zero ou no intervalo ( 0 ; 10 ] e o percentual de alunos nos intervalos subsequentes foi, \textbf{grosso} modo, decrescente.
Este aspecto do histograma é confirmado pelo coeficiente de assimetria ( 1,08 ), positivo e maior do que 1.
Para todas as regiões o coeficiente de assimetria também é positivo e em torno de 1,sendo maior na região Sudeste ( 1,27 ) e menor na região Centro-Oeste ( 0,77 ). 3.3.2.4 Comentários sobre a correção das respostas à Questão Discursiva 4 O ponto central da questão recai sobre a teoria estética, em especial, sobre as reflexões acerca do o problema da imagem no contexto da sociedade contemporânea.
Assim, a partir do enunciado o Ensino de Licenciatura em Artes Visuais : imagem, espetáculo e poder, o estudante deveria redigir um texto sobre esta temática a partir da leitura de um trecho do livro Muito além do espetáculo, organizado por Adauto Novaes.
Ali, são expostas ideias sobre a teoria de Guy Debord acerca da sociedade do espetáculo, em particular, a ideia de que, na sociedade capitalista contemporânea, o espetáculo – enquanto uma relação social entre indivíduos mediada por imagens – passa a determinar as vivências do indivíduo, em detrimento da realidade concreta, de tal forma que a ilusão, a aparência, a representação e o simulacro adquirem o estatuto de realidade.
A questão aponta, assim, para aspectos candentes da sociedade contemporânea, colocando em evidência o papel do arte-educador neste contexto.
O texto afirma que um dos desafios do ensino das Licenciatura em Artes Visuais é abordar criticamente as visualidades.
Coerente com esta abordagem, a questão discursiva se divide em duas partes : a ) A discussão sobre o papel dos diferentes atores sociais b ) A escolha de temas a serem trabalhados no ensino de Licenciatura em Artes Visuais.
Trata -se, portanto, de uma questão que pode gerar diversos debates, numa escala muito ampla, pois, num contexto histórico em que os limites entre realidade e ficção se atenuam cada vez mais, as imagens adquirem um valor excepcional, enquanto instrumentos passíveis de manipulação por parte dos mais diversos atores sociais ( personalidades, celebridades, políticos, artistas, profissionais do marketing, profissionais da comunicação, entre outros ) nos diferentes meios de comunicação de massa, principalmente, na televisão e na internet.
O caráter indissociável existente entre a produção das imagens enquanto instrumento de consumo e as diversas formas de poder torna a questão bastante polêmica, pois aponta para as formas de leitura das imagens, para a sua análise crítica, para formas de resistência ao consumo de imagens, entre outras.
Em suma, a questão traz o problema do papel do arte-educador da área de Licenciatura em Artes Visuais no contexto da sociedade de consumo ou sociedade do espetáculo.
Apesar da reelaboração do padrão de resposta, tendo em vista a realização de uma correção que pudesse avaliar o potencial do estudante de modo equânime, observamos que o desempenho geral foi considerado bastante abaixo do nível desejado, segundo os padrões de avaliação.
O número de respostas em branco foi bastante elevado.
No que diz respeito às respostas válidas, houve uma predominância de respostas fracas e insuficientes, isto é, que não atingem ao objetivo proposto.
Foram muito poucos os avaliados que obtiveram um desempenho acima do mediano ou com ótimo desempenho.
A Banca precisou adotar, inevitavelmente, o critério das notas relativas e das notas absolutas.
Dentro de um quadro onde o desempenho geral foi muito pouco significativo – talvez devido à própria complexidade da questão – algumas respostas obtiveram nota máxima, embora não tenham sido excelentes se consideradas isoladamente.
Nas respostas que receberam alta pontuação, percebeu -se que o estudante fez referências não somente às imagens no contexto da produção artística, mas também no contexto da cultura visual, mencionando o modo como as diversas mídias produzem imagens e como elas devem ser lidas e avaliadas.
Em alguns casos, os estudantes fizeram referências a pensadores e teóricos relacionados ao tema.
A partir desta abordagem inicial, que seriam de ordem teórica e conceitual, os alunos conceberam os papéis dos professores e alunos, a partir de propostas de ensino, tal como pede o enunciado.
Em muitas respostas não houve uma separação clara entre a parte “ a ” e a parte “ b ” da questão, pois ambas se mostram integradas de modo bastante coerente.
Outro aspecto considerado relevante para a avaliação encontra -se no próprio posicionamento crítico e ideológico do estudante, já que a questão envolve uma visão crítica sobre a sociedade capitalista.
Neste sentido, mesmo uma resposta que não pôde ser considerada excelente do ponto de vista técnico, mereceu nota elevada devido ao grau de compreensão do tema em termos sociais e políticos.
De forma comparativa, em respostas consideradas medianas foi observado que o estudante não chegou a aprofundar a questão de um ponto de vista conceitual ou que revelou certa inconsistência ou ingenuidade na conceituação – conforme ocorre com frequência com a ideia de “ espetáculo ” e de “ atores sociais ” – e que desenvolveu sua argumentação de maneira pouco precisa, embora fosse possível reconhecer uma compreensão geral do sentido da questão.
Em alguns casos, observou -se que o estudante aproveitou mal o espaço para a resposta, repetindo argumentos do enunciado da questão.
Ocorreu, também, de não haver acréscimo significativo ao que foi solicitado, com desvio do foco da questão.
Outro fator digno de nota encontra -se no fato de, em alguns casos, o estudante expor razoavelmente seus argumentos em relação à primeira parte da questão, mas não elaborar nada em relação à segunda parte.
As respostas que receberam baixa pontuação acentuaram os traços negativos das respostas medianas.
Assim, foram observados elementos tais como : a ) desvio do foco da questão ; b ) inconsistência da argumentação ; c ) abordagem generalizada do assunto, verificando -se que o estudante se limita a dizer o que seria óbvio, como por exemplo, “ a aproximação das aulas à realidade social e cultural do aluno ” ; d ) \textbf{reprodução} do enunciado.
No entanto, o que chamou mais a atenção nesse tipo de resposta foi a incompreensão conceitual.
Este fato ocorreu, em particular, com os conceitos de “ espetáculo ” e de “ atores sociais ”.
Enquanto os últimos foram confundidos com atores profissionais que trabalham em emissoras de TV, o primeiro tendeu a ser visto como um produto cultural, como um espetáculo teatral, sobretudo.
Assim, as respostas que manifestavam esta interpretação do tema foram rebaixadas em seu conceito ou obtiveram nota zero, o que ocorreu, em particular, com textos muito sucintos, que não acrescentavam nada ao enunciado, ou com respostas sem qualquer nexo ou relação com o tema, beirando o absurdo.
Neste sentido, as respostas que tiveram nota zero foram as consideradas como totalmente inadequadas para um estudante que se encontra no último ano de sua formação. 3.3.2.5 Análise da Questão Discursiva 5 do Componente de Conhecimento Específico A Tabela 3.14 contém as informações relativas à questão 5 do conjunto do componente de Conhecimento Específico.
A nota média dos estudantes de todo o Brasil foi 36,9.
A maior média foi registrada na região Sul ( 40,1 ), enquanto a menor média foi registrada na região Sudeste ( 34,5 ).
Quanto à variabilidade das notas, o desvio padrão dos alunos do Brasil, como um todo, foi 25,0.
Enquanto o maior desvio foi encontrado na região Nordeste ( 28,3 ), o menor foi encontrado na região Sul ( 22,9 ).
A nota máxima 100,0 foi alcançada nas regiões Norte, Sudeste e Sul, enquanto nas regiões Nordeste e Centro-Oeste, a nota máxima foi 95,0.
A nota mínima foi zero para todas as regiões do Brasil.
A mediana dos alunos do Brasil como um todo e das regiões Norte, Sul e Centro-Oeste foi 45,0.
Nas demais regiões as medianas foram : 50,0 no Nordeste e 40,0 no Sudeste.
O Gráfico 3.21 representa a distribuição das notas da questão discursiva 5 no componente de Conhecimento Específico.
Destaca -se o grande número de estudantes que deixaram a questão 5 em branco.
Considerando -se apenas os alunos que responderam a questão, a distribuição apresenta uma moda no intervalo ( 50 ; 60 ] e possui assimetria negativa ( coeficiente igual a – 0,30 ).
Nota -se um forte decrescimento na porcentagem dos alunos que obtiveram notas acima do intervalo modal. 3.3.2.6 Comentários sobre a correção das respostas à Questão Discursiva 5 A Questão 5 trata dos currículos praticados em Licenciatura em Artes Visuais e, no enunciado, pede que o estudante redija “ um texto dissertativo acerca do seguinte tema : Currículos praticados no ensino das Licenciatura em Artes Visuais ; a responsabilidade do professor perante as escolhas dos conteúdos a serem desenvolvidos ”.
Partindo do princípio de que a Arte passou a ser considerada como Área do Conhecimento e não mais como atividade, a questão solicitou ao estudante, de um lado, que elaborasse uma reflexão sobre a “ arte como área do conhecimento ou como atividade ” e, de outro, que considerasse “ o papel do professor como mediador do conhecimento artístico e cultural ”.
Podemos afirmar que se trata de uma questão fundamental, já que envolve o conhecimento acerca do próprio desenvolvimento histórico do ensino das Licenciatura em Artes Visuais no Brasil, desenvolvimento este que tem um marco nas reformas instituídas pela Lei nº 9.394/96 ( Lei de Diretrizes e Bases da Educação Nacional - LDB ), consideradas de grande impacto na legislação educacional brasileira, na medida em que priorizam a especificidade das linguagens artísticas e não mais a generalidade, tal como ocorria na legislação anterior.
Tais mudanças estabelecem um novo papel para o educador de Licenciatura em Artes Visuais, por compreenderem que a própria visualidade se mostra como um campo específico do conhecimento, envolvendo uma teoria da percepção e da sensibilidade, teorias da criação e da recepção, assim como teorias acerca do desenvolvimento histórico e estético das artes, enquanto elementos formadores de uma cultura visual.
Neste sentido, as novas Diretrizes almejam um profissional da arte-educação que seja apto à pesquisa, à produção ou criação, à crítica e ao ensino das Licenciatura em Artes Visuais.
Em suma, a questão lida com o conhecimento do estudante acerca das Diretrizes Curriculares Nacionais para os cursos de Artes, especialmente, o de Licenciatura em Artes Visuais.
Nas respostas elaboradas pelos estudantes, houve predominância das classificadas como “ medianas ” e “ insuficientes ”.
As respostas que receberam notas mais elevadas ( acima do intervalo ( 60 ; 70 ] ) foram identificadas com maior frequência na questão 5, se comparadas às duas outras de Conhecimento Específico.
A despeito desse resultado, sabendo -se que a Arte é considerada, concomitantemente, tanto uma atividade humana como uma esfera do conhecimento, constatou -se que a grande parte dos estudantes avaliados, no âmbito da questão 5, não conseguiu estabelecer, com sucesso, ou desejável clareza, uma prospecção rentável da inter-relação entre as esferas da habilitação/formação profissional e a área do conhecimento, como proposto pelos \textbf{Parâmetros} Curriculares Nacionais e pela LDB.
Tal fato se verifica em ambas as vertentes e habilitações consultadas no Exame, ou seja : 1.
Tanto em nível do Bacharelado, que prioriza o enfoque na formação do artista profissional, do seu desenvolvimento criativo e de suas habilidades específicas nas linguagens visuais ; e 2.
Em nível da Licenciatura, ou no âmbito do ensino de Licenciatura em Artes Visuais, que articulando as áreas da Arte e da Educação, visa a buscar que a formação do professor de Licenciatura em Artes Visuais se volte ao atendimento dos interesses do Ensino Fundamental e ao acolhimento das demandas inerentes ao Ensino Médio.
Sob essa perspectiva, na correção realizada pela Banca, foram observados vários comentários de discentes que, por estarem terminando o Bacharelado, não se sentiam capazes ou “ qualificados ” para responderem à questão 5.
Tal observância foi concretamente realizada porque houve várias declarações registradas a esse respeito.
Entretanto, para a grande maioria que respondeu a esta questão, foi notória uma expressiva incompreensão sobre o que realmente significa, em termos curriculares e \textbf{pedagógicos}, e nos aspectos que lhe são decorrentes, em nível dos seus objetivos específicos, a mudança de paradigma preconizada pela LDB e pelos \textbf{Parâmetros} Curriculares Nacionais, que consideraram explicitamente “ Artes ” como uma área de conhecimento.
Para expressiva parcela das respostas auferidas, percebeu -se, além de uma notória incompreensão sobre o que efetivamente significa “ área de conhecimento ”, um evidente desconhecimento sobre o significado, seja em que âmbito se utilize ou se considere, do termo “ cultura ”.
Além disso, uma significativa parcela de participantes se utilizou de um específico período do próprio enunciado da questão.
Para as questões respondidas de forma avaliada apenas como mediana, percebeu -se que o estudante identificou a importância da função do professor de Licenciatura em Artes Visuais, e a abordou de forma aligeirada, não discorrendo, com profundidade, sobre como efetivamente essa abordagem deva e pode efetivamente se dar.
Em nível de senso comum, percebeu -se na maioria das respostas concernentes à questão discursiva 5 da prova de Licenciatura em Artes Visuais que : conseguiu estabelecer, com sucesso, ou desejável clareza, uma prospecção rentável da inter-relação entre as esferas da habilitação/formação profissional e a área do conhecimento, como proposto pelos \textbf{Parâmetros} Curriculares Nacionais e pela LDB.
Tal fato se verifica em ambas as vertentes e habilitações consultadas no Exame, ou seja : 1.
Tanto em nível do Bacharelado, que prioriza o enfoque na formação do artista profissional, do seu desenvolvimento criativo e de suas habilidades específicas nas linguagens visuais ; e 2.
Em nível da Licenciatura, ou no âmbito do ensino de Licenciatura em Artes Visuais, que articulando as áreas da Arte e da Educação, visa a buscar que a formação do professor de Licenciatura em Artes Visuais se volte ao atendimento dos interesses do Ensino Fundamental e ao acolhimento das demandas inerentes ao Ensino Médio.
Sob essa perspectiva, na correção realizada pela Banca, foram observados vários comentários de discentes que, por estarem terminando o Bacharelado, não se sentiam capazes ou “ qualificados ” para responderem à questão 5.
Tal observância foi concretamente realizada porque houve várias declarações registradas a esse respeito.
Entretanto, para a grande maioria que respondeu a esta questão, foi notória uma expressiva incompreensão sobre o que realmente significa, em termos curriculares e \textbf{pedagógicos}, e nos aspectos que lhe são decorrentes, em nível dos seus objetivos específicos, a mudança de paradigma preconizada pela LDB e pelos \textbf{Parâmetros} Curriculares Nacionais, que consideraram explicitamente “ Artes ” como uma área de conhecimento.
Para expressiva parcela das respostas auferidas, percebeu -se, além de uma notória incompreensão sobre o que efetivamente significa “ área de conhecimento ”, um evidente desconhecimento sobre o significado, seja em que âmbito se utilize ou se considere, do termo “ cultura ”.
Além disso, uma significativa parcela de participantes se utilizou de um específico período do próprio enunciado da questão.
Para as questões respondidas de forma avaliada apenas como mediana, percebeu -se que o estudante identificou a importância da função do professor de Licenciatura em Artes Visuais, e a abordou de forma aligeirada, não discorrendo, com profundidade, sobre como efetivamente essa abordagem deva e pode efetivamente se dar.
Em nível de senso comum, percebeu -se na maioria das respostas concernentes à questão discursiva 5 da prova de Licenciatura em Artes Visuais que :  foi utilizada a ideia da \textbf{reprodução} pela \textbf{reprodução} – o desenhar pelo desenhar ou a atividade prática pela atividade prática – não foi reconhecível, em nível de senso comum, nem muito menos desejável como um bom procedimento \textbf{pedagógico} na condução da disciplina pelos docentes avaliados ;  houve significativa rejeição da função da disciplina Artes, sendo vista como mero instrumento ao atendimento de demandas observáveis a partir de datas e eventos comemorativos por solicitação das unidades educacionais ;  foi identificada uma manifestação de inferioridade em relação aos docentes dos componentes curriculares considerados “ mais valorizados ” tradicionalmente.
Neste sentido, valorizar a disciplina Artes dá -se mais em nível do reconhecimento recíproco dessa importância para os demais docentes do que efetivamente na construção de procedimentos e condutas que valorizem o ensino das Artes em si ;  o multiculturalismo é uma realidade que deve ser percebida, reconhecida e incorporada nas propostas \textbf{pedagógicas}, assim como o respeito às diversidades culturais que lhe são naturalmente decorrentes ;  a realidade cultural do educando deve ser respeitada e, considere -se aqui, especificamente, a utilização do lastro teórico proporcionado pela grande referência recorrida pelos estudantes avaliados advinda dos pressupostos na Pedagogia do Oprimido, de Paulo Freire, em que os referenciais do educando são determinantes à edificação de uma pedagogia realmente rentável e com sentido ;  as ideias observáveis a partir do pensamento da arte-educadora Ana Mae Barbosa – a grande referência específica utilizada da área para os alunos avaliados – e, essencialmente, sobre a sua teoria da Abordagem Triangular ( contextualização histórica, fazer artístico e apreciação artística ).
Esses itens apareceram na grande maioria das respostas válidas, individualmente ou articulados entre si, consideradas como medianas ou de alto nível de qualidade.
No entanto, para as avaliações que tinham um nível aquém do esperado, os itens citados também ocorreram, mas com insipiente qualidade de aprofundamento, ausência de articulação ou carência notória de desejável análise reflexiva.
Tornam -se mais um eco mecânico de um refrão acrítico e proferido sem sentido algum, como um mantra que não produz efeito que se valha a pena registrar. 3.3.3 Considerações Finais As três questões discursivas do Componente Específico lidavam com diferentes aspectos do tema Arte-Educação na Licenciatura em Artes Visuais.
A questão 3 propunha uma discussão acerca do fazer artístico ; a questão 4 inseria -se em discussões estéticas relacionadas à mídia e à cultura de massa, propondo um debate acerca das relações entre imagem e poder ; já a Questão 5 voltava -se para aspectos relativos à Legislação do ensino da Licenciatura em Artes Visuais.
Todas elas têm como ponto comum o fato de questionar o papel do professor de Licenciatura em Artes Visuais como mediador dos conhecimentos e como propositor de uma visão crítica.
No processo de correção, além dos dados sugeridos nos Padrões de Resposta, foram considerados os seguintes aspectos :  Quanto ao conteúdo : capacidade de o estudante responder ao que lhe era solicitado, desenvolvendo seus argumentos a partir da exposição de conhecimentos gerais sobre Licenciatura em Artes Visuais e ensino, assim como conhecimentos específicos ( técnicas artísticas, autores, conceitos ).  Comunicação e expressão : além de aspectos formais relativos à qualidade da escrita ( clareza, objetividade e concisão ), foi considerada, também, a visão pessoal do estudante, isto é, seus valores éticos, estéticos e sociais, visto que a área de ensino está intrinsecamente relacionada a estes valores.  Aspectos gramaticais : a forma como o aluno comunica suas ideias foi considerada como critério de julgamento, posto que constitui um indicativo de seus conhecimentos acerca da Língua Portuguesa e, evidentemente, refletem sua familiaridade com textos.
De um modo geral, em todas as questões discursivas do Componente Específico houve pouquíssima referência a artistas e a teóricos da Arte e da estética.
As poucas referências situam -se no quadro do Modernismo e das Vanguardas.
Este fato aponta para uma falta de recursos do ponto de vista teórico e conceitual, falta esta que se reflete na pouca qualidade reflexiva presente nas respostas.
Nas três questões, o que se notou foi uma predominante preocupação com a aplicação de uma metodologia de trabalho, a adoção de um saber de ordem operacional, sobretudo na Questão 3.
Esse modo de abordagem se torna problemático em função mesmo da complexidade da Arte – e da Educação – no contexto contemporâneo.

% matched lemmas: grosso, parâmetro, pedagógico, reprodução
\end{textsample}
